\noindent \textbf{Theorem.} Let $\mathcal{M}(S)$ denote the set of multiples of a subset $S \subseteq \mathbb{N}$, and let $\overline{d}(X)$ denote the upper asymptotic density of a set $X$. There exists a sequence $A : \mathbb{N} \to \mathbb{N}$ satisfying $\sum_{i \in \mathbb{N}} 1/A(i) = \infty$, such that for $\epsilon = 1/4$ and for all $k \ge 1$, we have $\overline{d}(\mathcal{M}(A + k)) < 1 - \epsilon$. 

\noindent \textbf{Proof.} We construct the sequence iteratively in finite blocks and show that for any arbitrary shift $k$, we can bound the upper density of the shifted sequence by analyzing three distinct cases, the third case relying on the construction being similar to the simple counterexample of Ruzsa which resolved the original question Erd\H{o}s \#26.
\\\\
Let $(q_j)_{j=1}^\infty$ be a strictly increasing sequence of primes satisfying $q_1 \ge 29$ (chosen large enough to cause small enough density in the infinite tail of the sequence). We define a strictly increasing sequence of indices (or ``jumps'') $(J_m)_{m=0}^\infty$ recursively, starting with $J_0 = 0$. Suppose $J_m$ has been defined, and the sequence $A(i)$ has been constructed for all $i < J_m$. We define the parameters for the $m$-th block as follows.
\\\\
First, define the step size $P_m = \prod_{j = 1}^{10J_m} q_j$. Next, by the Chinese Remainder Theorem, we can choose a base size $R_m$ satisfying the following system of congruences
$$R_m \equiv -j \pmod{q_j} \quad \text{for } 1 \leq j \leq 10J_m$$
as well as the strict monotonicity condition $R_m > A(J_m - 1)$ (where we set $A(-1) = 0$), by adding enough multiples of $P_m$.
\\\\
With $R_m$ and $P_m$ fixed, the series $\sum_{x=1}^\infty \frac{1}{R_m + xP_m}$ diverges. Thus, we can compute a block length $L_m \ge 1$ such that the partial sum satisfies:
$$0.1 \leq \sum_{x = 1}^{L_m} \frac{1}{R_m + xP_m} \leq 0.2$$
We then set $J_{m + 1} = J_m + L_m$. Now that the bounds of the block are properly defined, for indices $i \in [J_m, J_{m+1})$, we set:
$$A(i) = R_m + (i - J_m + 1)P_m$$
Since each block contributes at least $0.1$ to the total harmonic sum, the overall sum $\sum 1/A(i)$ diverges. Pick an arbitrary shift $k \ge 1$; we must now show that $\overline{d}(\mathcal{M}(A+k))$ stays bounded below $3/4$. Let $m_0$ be the unique integer satisfying:
$$10J_{m_0} \leq k < 10J_{m_0 + 1}$$
We partition the shifted sequence $A+k$ into three regimes and bound the upper density of the multiples generated by each:
\begin{itemize}
    \item (Case 1: $i < J_{m_0}$) The upper density of the multiples generated by this finite set is bounded by the sum of their reciprocals. Because $A(i) > 0$, we have $A(i) + k > k$, yielding:
    $$\overline{d}(\mathcal{M}(\{A(i) + k \mid i < J_{m_0}\})) \leq \sum_{i = 0}^{J_{m_0} - 1} \frac{1}{A(i) + k} < \sum_{i = 0}^{J_{m_0} - 1} \frac{1}{k} = \frac{J_{m_0}}{k} \leq 0.1$$
    \item (Case 2: $J_{m_0} \leq i < J_{m_0+1}$) Shifting by a positive $k$ strictly decreases the reciprocals, so the upper density of the multiples from this critical block satisfies:
    $$\overline{d}(\mathcal{M}(\{A(i) + k \mid J_{m_0} \le i < J_{m_0+1}\})) \le \sum_{i=J_{m_0}}^{J_{m_0+1}-1} \frac{1}{A(i) + k} < \sum_{i=J_{m_0}}^{J_{m_0+1}-1} \frac{1}{A(i)} \le 0.2$$
    \item (Case 3: $i \geq J_{m_0 + 1}$) Consider any element $A(i)$ where $i$ belongs to some block $m \ge m_0 + 1$, meaning $10 J_m \ge 10 J_{m_0+1} > k$. Recall the construction of the $m$-th block:
    $$A(i) + k = (R_m + k) + (i - J_m + 1)P_m$$
    Because $10 J_m \ge k$, we have $q_k \mid P_m$. By construction, we also have $q_k \mid (R_m + k)$. Consequently, $q_k \mid (A(i) + k)$ for every element in this infinite tail. The set of multiples generated by Case 3 is therefore entirely contained within the multiples of the prime $q_k$. Expressing this with the upper density notation:
    $$\overline{d}(\mathcal{M}(\{A(i) + k \mid i \ge J_{m_0+1}\})) \le \overline{d}(\mathcal{M}(\{q_k\})) = \frac{1}{q_k} \le \frac{1}{29}$$
\end{itemize}
By the union bound, the total upper density of the multiples of $A+k$ is bounded above by:
$$0.1 + 0.2 + \frac{1}{29} \approx 0.334 < \frac{3}{4} = 1 - \epsilon.$$ \hfill $\Box$
